\section{Data, Classification, and Samples}
\label{sec:data}

\subsection{Procurement and litigation data}
\label{subsec:data_sources}

The data cover S\~ao Paulo pharmaceutical procurement through BEC during \BPsampleStartYear--\BPsampleEndYear. The empirical file contains \BPnPOIfull{} purchase-offer-item observations. BEC records item identifiers, buyers, suppliers, reference prices, negotiated prices, quantities, participating firms, winners, procurement modality, and tender outcomes. We link these procurement records to SES/SP litigation and administrative-process markers to classify purchases into ordinary, administrative urgent, and litigated urgent regimes.

The classifier operates upstream at the purchase-order/tender-notice
level, where legal status is observed. After regime labels are linked
to BEC item records, outcomes are measured at the
purchase-offer-item level. Price regressions use accepted winning
bids; other outcomes use the broader purchase-offer-item or
tender-level sample specified in each table. This preserves the level
of the legal mandate while using the transaction level at which
prices, quantities, and winners are realized.

\subsection{Classification of purchase regimes}
\label{subsec:classification}

Purchase regimes are classified using structured links and
tender-notice text. The classifier combines machine-learning
predictions, JSON fields, regex markers, and position-based checks to
separate ordinary, administrative urgent, and litigated urgent
purchases.

Validation is conducted at the purchase-order/tender-notice level.
Against \BPregexNGroundTruth{} ground-truth purchase orders, the
classifier reaches \BPregexExactMatchPct{} exact agreement. Per-class
F1 is \BPregexFoneJud{} for judicial purchases and \BPregexFoneAdm{}
for administrative purchases; the macro-F1 over those two urgent
classes is \BPregexFone{}. Online Appendix A reports the validation
and sample-construction tables.

Remaining misclassification would attenuate regime contrasts unless it
is systematically correlated with residual prices within controls. The
main concern is differential error between urgent regimes among
high-cost items; the validation evidence reduces, but cannot eliminate,
that concern.

\subsection{Samples and identifying variation}
\label{subsec:samples}

The paper uses several samples because each empirical comparison asks
a different question. The analysis sample supports
ordinary-versus-urgent comparisons. The winners-only price sample
contains \BPnegNobs{} accepted winning bids and is used when negotiated
price is the outcome.

Within urgent procurement, the urgent panel contains administrative
and litigated winning bids and supports the under-the-gun comparison.
The firm-buyer-item triple sample is narrower by design. It requires
the same supplier to sell the same item to the same buyer under both
urgent regimes, so it identifies within-supplier pricing. It does not
measure supplier-set reallocation, which requires the winner-switching
evidence in Section~\ref{sec:results}.

The samples narrow as the question narrows. The
ordinary-versus-urgent comparison supports the broad urgent-demand
claim. The urgent panel focuses on sanction exposure among purchases
already under urgent execution. The triple sample conditions on the
same supplier; that narrowness is the source of the pricing test, not
an estimate of total costs.

\subsection{Descriptive facts}
\label{subsec:descriptives}

The descriptive facts show why raw administrative-versus-litigated means are not a design. Administrative urgent purchases are larger than litigated urgent purchases, while litigated purchases have higher average log negotiated prices. These differences combine selection, scale, and sanction exposure, so we do not interpret raw means causally. The design instead bounds selection, decomposes scale, and tests same-firm pricing and supplier-set reallocation directly.

The descriptive patterns nevertheless motivate the mechanism analysis:
larger administrative orders suggest more aggregation within the
urgent tier, while higher raw litigated prices show the direction of
the problem but not its cause.
